\documentclass[10pt]{article}

\usepackage[a4paper, total={6.5in, 9.5in}]{geometry}
\usepackage[dvipsnames,svgnames]{xcolor}
\usepackage[shortlabels]{enumitem}
\usepackage[framemethod=TikZ]{mdframed}
\usepackage{amsmath,amssymb,amsthm}
\usepackage[colorlinks]{hyperref}
\usepackage{tikz-cd}

\hypersetup{
	colorlinks=true,
	linkcolor=blue,
	filecolor=blue,      
	urlcolor=blue,
	pdftitle={Obstruction theory notes}
}

\usepackage{mathtools}
\usepackage{thmtools}
\usepackage[textsize=scriptsize]{todonotes}

\renewcommand{\tilde}{\widetilde}
\renewcommand{\hat}{\widehat}
\renewcommand{\setminus}{\!-\!}
\providecommand{\ol}{\overline}
\providecommand{\eps}{\varepsilon}
\providecommand{\half}{\frac{1}{2}}
\providecommand{\dang}{\measuredangle} %% Directed angle
\providecommand{\dd}{\mathrm d}
\providecommand{\CC}{\mathbb C}
\providecommand{\FF}{\mathbb F}
\providecommand{\HH}{\mathbb H}
\providecommand{\NN}{\mathbb N}
\providecommand{\QQ}{\mathbb Q}
\providecommand{\RR}{\mathbb R}
\providecommand{\ZZ}{\mathbb Z}
\providecommand{\RP}{\mathbb RP}
\providecommand{\CP}{\mathbb CP}
\providecommand{\dg}{^\circ}
\providecommand{\ind}[1]{\mathbf 1_{#1}}
\providecommand{\ii}{\item}
\providecommand{\alert}[1]{{\sffamily\textbf{\textcolor{magenta}{#1}}}}
\providecommand{\opname}{\operatorname}
\providecommand{\li}{\opname{li}}
\providecommand{\ls}{\opname{ls}}
\providecommand{\supp}{\opname{supp}}
\providecommand{\im}{\opname{im}}
\providecommand{\Ext}{\opname{Ext}}
\providecommand{\Tor}{\opname{Tor}}
\providecommand{\inv}{^{-1}}
\providecommand{\setdiff}{\smallsetminus}
\providecommand{\mb}[1]{\mathbf{#1}}
\providecommand{\abs}[1]{\left\lvert{#1}\right\rvert}
\providecommand{\norm}[1]{\left\|{#1}\right\|}
\providecommand{\dif}{\;d}
\providecommand{\id}{\operatorname{id}}

\newcommand{\expow}{\Lambda}
\newcommand{\tensor}{\otimes}
\newcommand{\Hom}{\opname{Hom}}

\reversemarginpar
\setlength{\marginparwidth}{1in}

\mdfdefinestyle{mdbluebox}{roundcorner=10pt,innerbottommargin=9pt,
	linecolor=blue,backgroundcolor=TealBlue!5,}
\declaretheoremstyle[headfont=\sffamily\bfseries\color{MidnightBlue},
mdframed={style=mdbluebox},]{thmbluebox}

\mdfdefinestyle{mdredbox}{frametitlefont=\bfseries,innerbottommargin=8pt,
	nobreak=true,backgroundcolor=Salmon!5,linecolor=RawSienna,}
\declaretheoremstyle[headfont=\bfseries\color{RawSienna},
mdframed={style=mdredbox},headpunct={\\[3pt]},postheadspace=0pt,]{thmredbox}

\mdfdefinestyle{mdgreenbox}{linecolor=ForestGreen,backgroundcolor=ForestGreen!5,
	linewidth=2pt,rightline=false,leftline=true,topline=false,bottomline=false,}
\declaretheoremstyle[headfont=\bfseries\sffamily\color{ForestGreen!70!black},
mdframed={style=mdgreenbox},headpunct={ --- },]{thmgreenbox}

\mdfdefinestyle{mdorangebox}{linecolor=RedOrange,backgroundcolor=RedOrange!5,
	linewidth=2pt,rightline=false,leftline=true,topline=false,bottomline=false,}
\declaretheoremstyle[headfont=\bfseries\sffamily\color{RedOrange!70!black},
mdframed={style=mdorangebox},headpunct={ --- },]{thmorangebox}

\mdfdefinestyle{mdblackbox}{linecolor=black,backgroundcolor=RedViolet!5!gray!5,
	linewidth=3pt,rightline=false,leftline=true,topline=false,bottomline=false,}
\declaretheoremstyle[mdframed={style=mdblackbox}]{thmblackbox}

\declaretheoremstyle[spaceabove=6pt,spacebelow=6pt]{boldhead}

\declaretheorem[style=thmbluebox,name=Theorem,numberwithin=section]{theorem}

\declaretheorem[style=boldhead,name=Example,sibling=theorem]{example}

\declaretheorem[style=thmbluebox,name=Corollary,sibling=theorem]{corollary}
\declaretheorem[style=thmbluebox,name=Proposition,sibling=theorem]{prop}
\declaretheorem[style=thmbluebox,name=Lemma,sibling=theorem]{lemma}
\declaretheorem[style=thmbluebox,name=Theorem,numbered=no]{theorem*}
\declaretheorem[style=thmbluebox,name=Lemma,numbered=no]{lemma*}
\declaretheorem[style=thmgreenbox,name=Claim,numberwithin=theorem]{claim}
\declaretheorem[style=thmgreenbox,name=Claim,numbered=no]{claim*}
\declaretheorem[style=boldhead,name=Remark,sibling=theorem]{remark}
\declaretheorem[style=boldhead,name=Remark,numbered=no]{remark*}
\declaretheorem[style=thmgreenbox,name=Definition,sibling=theorem]{definition}
\declaretheorem[style=thmgreenbox,name=Definition,numbered=no]{definition*}
\declaretheorem[style=thmorangebox,name=Warning,sibling=theorem]{warning}
\declaretheorem[style=thmorangebox,name=Warning,numbered=no]{warning*}
\declaretheorem[style=thmorangebox,name=Note,numbered=no]{note}
\declaretheorem[style=thmblackbox,name=Exercise,sibling=theorem]{exercise}
\declaretheorem[style=thmblackbox,name=Exercise,numbered=no]{exercise*}

\newenvironment{soln}{\begin{proof}[Solution]}{\end{proof}}
\newenvironment{walkthrough}{\noindent \textbf{Walkthrough.}}{\medskip}
\newlist{walk}{enumerate}{3}
\setlist[walk]{label=\bfseries (\alph*)}

\providecommand{\pow}{\operatorname{Pow}}
%\providecommand{\vocab}[1]{{{\sffamily\textolor{ForestGreen}{#1}}}}
\providecommand{\vocab}[1]{{{\sffamily\textit{#1}}}}
\providecommand{\lcm}{\operatorname{lcm}}
\providecommand{\ord}{\operatorname{ord}}
\providecommand{\Int}{\operatorname{Int}}
\providecommand{\rk}{\operatorname{rk}}
\providecommand{\Mat}{\operatorname{Mat}}

\usepackage{asymptote}
\begin{asydef}
	defaultpen(fontsize(10pt));
	unitsize(1cm);
	size(8cm); // set a reasonable default
	usepackage("amsmath");
	usepackage("amssymb");
	settings.tex="pdflatex";
	settings.outformat="pdf";
	// Replacement for olympiad+cse5 which is not standard
	import geometry;
	// recalibrate fill and filldraw for conics
	void filldraw(picture pic = currentpicture, conic g, pen fillpen=defaultpen, pen drawpen=defaultpen)
	{ filldraw(pic, (path) g, fillpen, drawpen); }
	void fill(picture pic = currentpicture, conic g, pen p=defaultpen)
	{ filldraw(pic, (path) g, p); }
	// some geometry
	pair foot(pair P, pair A, pair B) { return foot(triangle(A,B,P).VC); }
	pair orthocenter(pair A, pair B, pair C) { return orthocentercenter(A,B,C); }
	pair centroid(pair A, pair B, pair C) { return (A+B+C)/3; }
	// cse5 abbreviations
	path CP(pair P, pair A) { return circle(P, abs(A-P)); }
	path CR(pair P, real r) { return circle(P, r); }
	pair IP(path p, path q) { return intersectionpoints(p,q)[0]; }
	pair OP(path p, path q) { return intersectionpoints(p,q)[1]; }
	path Line(pair A, pair B, real a=0.6, real b=a) { return (a*(A-B)+A)--(b*(B-A)+B); }
	// cse5 more useful functions
	picture CC() {
		picture p=rotate(0)*currentpicture;
		currentpicture.erase();
		return p;
	}
	pair MP(Label s, pair A, pair B = plain.S, pen p = defaultpen) {
		Label L = s;
		L.s = "$"+s.s+"$";
		label(L, A, B, p);
		return A;
	}
	pair Drawing(Label s = "", pair A, pair B = plain.S, pen p = defaultpen) {
		dot(MP(s, A, B, p), p);
		return A;
	}
	path Drawing(path g, pen p = defaultpen, arrowbar ar = None) {
		draw(g, p, ar);
		return g;
	}
\end{asydef}


\title{\vspace{-2.0cm}The Thom Isomorphism and Characteristic Classes}
\author{\large Aditya Pahuja}
\date{\large \today}
\begin{document}

\maketitle
\tableofcontents
\pagebreak

\section{Homology and cohomology}

\section{Smooth vector bundles}

\subsection{Definitions and basic properties}

\subsection{Oriented vector bundles}


Let $\pi\colon E\to B$ be an oriented rank-$n$ vector bundle.
The Thom isomorphism relates the cohomology
$H^k(B; R)$ to the relative cohomology of $H^{k+n}(E,E^\circ;  R)$
where $E\dg$ is the complement of the zero section of $E$
and $ R$ is an (associative) unital ring.
Note that for each fiber $F = \pi\inv(b)$, $b\in B$,
with $F\dg = F\cap E\dg$ being the nonzero elements of $F$,
the relative cohomology $H^*(F,F\dg;  R)$ is equal to
\begin{equation*}
    H^*(F,F\dg;  R) = \begin{cases}
	 R & \text{if }* = n\\
	0 & \text{otherwise;}
    \end{cases}
\end{equation*}
one can see this by considering the long exact sequence of the pair $(F,F\dg)$.
The orientation on the vector bundle then confers a
canonical generator $u_{F,\ZZ}$ upon each
fiberwise cohomology group $H^n(F,F\dg;\ZZ)\cong\ZZ$.
Then, since there is a unique homomorphism $\ZZ\to R$
for any unital ring $ R$, we may define $u_{F, R}$
to be the image of $u_{F,\ZZ}$ under the map
$H^n(F,F\dg;\ZZ)\to H^n(F,F\dg; R)$;
thus the orientation of the vector bundle gives a canonical generator
for the fiberwise cohomologies regardless of the ring.

% TODO move most of the above to a separate preceding section
% TODO this is a bad way to do the canonical orientatino stuff;
% rephrase wrt  R-ortienability

\section{Thom isomorphism}
We follow Milnor and Stasheff's account of the Thom isomorphism theorem.

\begin{definition}[Thom class]
    \label{def:thom-class}
    A \vocab{Thom class} $u = u_R\in H^n(E,E\dg; R)$
    which restricts to $u_F = u_{F, R}$ for every fiber
    $F = \pi\inv(b)$ above a point $b\in B$.
\end{definition}

\begin{theorem}[Thom isomorphism]
    \label{thm:thom-isomorphism}
    Each $ R$-oriented vector bundle $\pi\colon E\to B$
    possesses a unique Thom class. Furthermore, there is an isomorphism
    \begin{align*}
	\Phi\colon H^k(E; R)&\to H^{k+n}(E,E\dg; R)\\
	y&\mapsto y\smile u.
    \end{align*}
\end{theorem}

Note that $E$ deformation retracts onto the zero section,
which is homeomorphic to $B$, by simply linearly contracting each fiber.
Thus the Thom isomorphism gives an isomorphism
$\Phi\circ\pi^*\colon H^k(B; R)\to H^k(E,E\dg; R)$.

Some geometric intuition for the Thom class:
Consider the space of differential $n$-forms on $E$
whose restriction to each fiber is compactly supported;
these are called \vocab{differential forms with compact vertical supports},
which form cohomology groups $H^*_{\text{cv}}(E)$.
A \vocab{Thom form} $\phi$ is an $n$-form on $E$
whose restriction to each fiber integrates to $1$
(thus represents the Thom class of $E$ when $ R = \RR$).
The Thom isomorphism theorem says that there is an isomorphism
$H_{\text{dR}}^k(B)\to H^{k+n}_{\text{cv}}(E)$
via wedging with $\phi$; the inverse of this map
is given by integrating a compactly vertically supported $(k+n)$-form
on each fiber to get a $k$-form on $B$.
% TODO elaborate?

% TODO talk about Thom spaces elsewhere
% An alternative perspective is to observe that
% the relative cohomology $H^k(E,E\dg; R)$ manifests
% as the cohomology of the \vocab{Thom space}, defined as follows.
% After equipping $E$ with a Euclidean structure,
% % TODO explain somewhere why this can always be done
% let $D(E)$ be the subspace of $E$ consisting of vectors
% with magnitude at most $1$ (hence is an $n$-disk bundle)
% and $S(E)$ the subspace consisting of vectors
% with magnitude equal to $1$ (hence is an $(n-1)$-sphere bundle).
% Then the Thom space $T(E) = D(E)/S(E)$ can be thought of
% as compactifying each fiber of $E$ to get an $n$-sphere bundle
% and then collapsing the zero section to a point.
% The Thom isomorphism theorem is a statement about
% the cohomology of the Thom space, since
% \begin{equation*}
%     H^k(E,E\dg; R) = H^k(D(E), S(E); R)
%     = \widetilde H^k(D(E)/S(E); R).
% \end{equation*}
% 
% \missingfigure{thom space (steal from F\&F maybe)}

The proof of Theorem \ref{thm:thom-isomorphism}
will proceed in roughly three steps:
first, we prove it for trivial vector bundles;
then, we prove it when $B$ can be covered by
finitely many open sets over which the given vector bundle is trivial;
and finally, we prove it in the general case
by taking a limit over compact sets.

\paragraph{The case of a trivial bundle.}


\paragraph{The case of a good finite cover of $B$.}
We will show that 
if that $B = B_1 \cup B_2$ is a cover of $B$ by two open sets
such that Theorem \ref{thm:thom-isomorphism} holds for
the restricted vector bundles
$\pi|_{\pi\inv(B_1)}\colon E_1 = \pi\inv(B_1)\to B_1$ and
$\pi|_{\pi\inv(B_2)}\colon E_2 = \pi\inv(B_2)\to B_2$,
then the result of the theorem also holds for $\pi\colon E\to B$.
This will imply the result in the case where $B$ is compact,
or more generally whenever $B$ can be covered by finitely many
trivializations of the vector bundle.

Let $u_1\in H^n(E_1,E_1\dg; R)$ and $u_2\in H^n(E_2,E_2\dg; R)$
be the Thom classes of the two restricted bundles.
The restrictions of $u_1$ and $u_2$ to $E_1\cap E_2 = E_\cap$
are both Thom classes for the restricted bundle over $B_1\cap B_2$,
so by uniqueness of the Thom class on $E_1$ and $E_2$,
$u_1|_{E_\cap} = u_2|_{E_\cap} = u_\cap$.
Therefore, there is a well-defined cohomology class $u$ on $E$
whose restrictions to $E_1$ and $E_2$ are respectively
$u_1$ and $u_2$. This is clearly a Thom class because it is a Thom class
fiberwise, and its uniqueness follows from the uniqueness of $u_1$ and $u_2$,
for any other cohomology class restricts to something different
from $u_1$ on $E_1$ or something different from $u_2$ on $E_2$.
To show that cupping with $u$ is an isomorphism,
apply the five lemma to the diagram below,
in which the rows are segments of Mayer-Vietoris sequences
(and the coefficient ring $R$ is omitted for neatness).
\begin{center}
    \small
    \begin{tikzcd}[column sep=1.1em]
	H^k(E_1)\oplus H^k(E_2)\arrow[r]
	\arrow[d, "(\smile u_1)\oplus (\smile u_2)"]
	& H^k(E_\cap)\arrow[r]
	\arrow[d, "\smile u_\cap"]
	& H^{k+1}(E)\arrow[r]
	\arrow[d, "\smile u"]
	& H^{k+1}(E_1)\oplus H^{k+1}(E_2)\arrow[r]
	\arrow[d, "(\smile u_1)\oplus (\smile u_2)"]
	& H^{k+1}(E_\cap)
	\arrow[d,"\smile u_\cap"]
	\\
	H^{k+n}(E_1,E_1\dg)\oplus H^{k+n}(E_2,E_2\dg)\arrow[r]
	& H^{\ell}(E_\cap, E_\cap\dg)\arrow[r]
	& H^{{\ell}+1}(E)\arrow[r]
	& H^{{\ell}+1}(E_1,E_1\dg)\oplus H^{{\ell}+1}(E_2,E_2\dg)\arrow[r]
	& H^{\ell+1}(E_\cap, E_\cap\dg)
    \end{tikzcd}
\end{center}
The commutativity of this diagram follows from
commutativity of the corresponding map between
short exact sequences of chain complexes. Specifically,
consider the diagram of chain complexes
\begin{center}
    \begin{tikzcd}
	0\arrow[r]
	& C^*(E)\arrow[r]
	\arrow[d,"\smile u"]
	& C^*(E_1)\oplus C^*(E_2)\arrow[r]
	\arrow[d, "(\smile u_1)\oplus (\smile u_2)"]
	& C^*(E_\cap)\arrow[r]
	\arrow[d, "\smile u_\cap"]
	& 0\\
	0\arrow[r]
	& C^{*+n}(E,E\dg)\arrow[r]
	& C^{*+n}(E_1,E_1\dg)\oplus C^{*+n}(E_2,E_2\dg)\arrow[r]
	& C^{*+n}(E_\cap,E_\cap\dg)\arrow[r]
	& 0
    \end{tikzcd}
\end{center}
and let
$i_k\colon E_k\to E$ and $j_k\colon E_\cap\to E_k$
be the inclusions for $k=1,2$. Then, the diagram commutes because
\begin{equation*}
    (i_1^\sharp\alpha\smile u_1, i_2^\sharp\alpha\smile u_2)
    = (i_1^\sharp(\alpha\smile u), i_2^\sharp(\alpha\smile u)),\quad
    (j_1^\sharp\beta - j_2^\sharp\beta)\smile u_\cap
    = j_1^\sharp(\beta\smile u_1) - j_2^\sharp(\beta\smile u_2),
\end{equation*}
where $f^\sharp$ is the cochain map induced by $f$.










\end{document}
